\documentclass[11pt]{article}
\usepackage[margin=2.3cm]{geometry}
\usepackage{amsmath,amssymb,amsthm}
\usepackage[T1]{fontenc}
\usepackage{enumitem}
\usepackage{booktabs}
\usepackage{hyperref}
\newcommand{\Norm}{\mathbf{Norm}}
\renewcommand{\Im}{\operatorname{Im}}
\setlength{\parskip}{4pt}\setlength{\parindent}{0pt}

\title{Review: \emph{What the weld's non-associativity is --- and is not:\\
a triangulation across four obstruction theories, and the idempotent normalizer that receives it}}
\author{Rick\\[2pt]\small To: Scot MacBeth (Kodamai)}
\date{2026-10-21}

\begin{document}
\maketitle
\thispagestyle{empty}
\begin{center}
\begin{tabular}{ll}
Reviewed: & MacBeth content \texttt{ac07a05} (pushed \texttt{24c88bc}), 10 pp., email UID 425\\
This review: & WIP commit: TBD\\
Check script: & \texttt{checks/2026-10-21-macbeth-triangulation-check.py} (+\,\texttt{.log})
  \\ & exact arithmetic, ${<}10$\,s
\end{tabular}
\end{center}

Locators are page numbers of the PDF, plus the number of the displayed item. You call the results
``Theorem A/B'' and ``Prop 5.1'' in the email. In the PDF they are Theorem 8.1, Theorem 8.2 and
Proposition 10.1, and I use the PDF numbers throughout.

\section*{Verdict}
\begin{itemize}[leftmargin=*]
\item \textbf{The algebra of Movement II is correct.} I checked these by hand and by machine:
  Lemma 5.1, Theorem 5.2, Lemma 6.3, Theorems 8.1 and 8.2, Lemma 9.2, the Lemma 9.3 witness
  (exact: $(f\#g)\#h=(\tfrac{8}{11},\tfrac{3}{11})$, $f\#(g\#h)=(\tfrac23,\tfrac13)$,
  $D_{gh}=\mathrm{diag}(\tfrac32,2)$), Proposition 10.1 and the Lean model of Remark 10.4. Your claim
  that no model exists at order ${\le}3$ is right. At order 4 the model is \emph{unique} up to
  isomorphism (C6, C10).
\item \textbf{The framing claims do not yet stand.} The note never names a map from the weld to
  $\Norm$ (M1). The general argument in Refutation I is false as stated (M2). Proposition 10.1 is
  vacuous under your own axiom (U), and it does not prove ``never element-level'' (M3). The positive
  home is, structurally, the classical transversal construction of right loops, so ``original
  ground'' cannot stand as written (M4).
\item \textbf{Grade.} \S\S7--9 and Prop.\,10.1 as an algebraic statement: \textbf{checked-sober}.
  The triangulation thesis, the ``element-level never'' clause and the originality: \textbf{sketched}.
\end{itemize}

\section*{Your two questions}
\textbf{Q1 (triangulation or list?).} As written it is a list with a good taxonomy. It is not yet a
triangulation. Three reasons.
(a) Movement~I holds the associative corner plus \emph{three} refutations (\S\S4--6). The fourth is in
Movement~II (\S10). The email's phrase ``Movement I's four kinds'' does not match the PDF.
(b) Rows 1 and 2 of the \S11 table probe the same object, the single transformation $k$, and
\S11 says they are ``the same fact read twice''. So at most three kinds are independent.
(c) The three refutations in Movement~I all die for \emph{one} reason. Each tests a property of $k$,
and $k$ is a unary datum: a 1-cell $DM\Rightarrow DM$ that never sees a triple. The defect is a
ternary datum, $Z_K=\{(f,g,h):K(f\cdot K(g\cdot h))\neq K(f\cdot g\cdot h)\}$. That single
reason is the real content. It is one theorem read three ways, not three independent bearings.
A genuine triangulation has to test each candidate home on the invariant itself (on $Z_K$, or on the
cocycle $\varphi(f,g)=D_{fg}$ of M4). It also needs, for each pair of homes, an example that one
home's stated reason kills and the other's does not. As things stand the cocycle row is not even
settled (M2, M4). My recommendation: retitle Movement~I as ``one reason, three readings''. That is
honest, and it is a stronger statement than a triangulation claim that cannot be cashed.

\textbf{Q2 (is Prop.\,10.1 decisive?).} It is correct. I find 1434 principal normalizers
($\kappa x$, $x\kappa$ and $\kappa x\kappa$ with $\kappa^2=\kappa$, over all monoids of order
${\le}4$) and none is non-associative (C8). But it is not decisive for the claim it is used for.
\begin{enumerate}[label=(\roman*),leftmargin=*]
\item \emph{Under (U) it is vacuous.} $K(1_X)=\kappa\cdot 1_X=\kappa$, so (U) forces $\kappa=1$ and
  $K=\mathrm{id}$. Machine check: no principal $K\neq\mathrm{id}$ satisfies (U) (C8). Dropping (U), it
  is true, but $N_K$ then has no units.
\item \emph{The real mechanism is closure, not elements.} If $\Im K$ is closed under $\cdot$, then
  for $f,g\in\Im K$ we get $f\#g=K(fg)=fg$, so $\#$ is the restriction of $\cdot$ and is associative.
  This one line covers $\kappa x$, $x\kappa$ and $\kappa x\kappa$ at once. Conversely, every handed
  normalizer found has a non-closed image (C6: 6/6 at order 4). Your own Remark 7.3 already says
  this: ``composing two normalized kernels de-normalizes''.
\item \emph{``Element vs operator'' is the wrong dichotomy.} Norm's $K$ is $x\mapsto D_x\cdot x$:
  left multiplication by an \emph{element} that depends on $x$. The group-transversal $K$ of M4 is
  the same, $g\mapsto h(g)^{-1}g$ with $h(g)\in H$. What Prop.\,10.1 rules out is a
  \emph{constant} multiplier. The decisive statement is: handed $\Rightarrow$ $\Im K$ is not closed.
\item The near-semiring translation $K=\kappa\cdot(-)$ is one formalization that you chose. A
  literally left-absorbing $\kappa$ ($\kappa a=\kappa$) gives a constant $K$ with
  $\Im K=\{\kappa\}$. So the refutation hits your formalization and not the near-semiring idea in
  general. Say which formalization is being refuted.
\end{enumerate}

\section*{Findings}
\subsection*{MUST}
\begin{enumerate}[label=\textbf{M\arabic*.},leftmargin=*]
\item \textbf{Title; p.\,1 l.\,``The weld is not associative in general''; p.\,2 Rem.\,1.3; p.\,3
  Rem.\,3.3; \S12.} There is no named map from the weld to $\Norm$, and the note contradicts itself on
  the weld. Rem.\,3.3 says the weld, once well defined, ``is a genuine comodule'' and lives in the
  associative regime. Def.\,1.1 and the abstract say the weld is non-associative. The note never
  defines which operation on welds fails to associate, or on which triple. Rem.\,1.3 calls $k$ a
  ``faithful stand-in'' and cites \S9 as computational justification. But \S9 computes $\Norm$, not the
  weld, so the justification is circular. The ``principle'' (both are idempotents measuring departure
  from a distributive law) is the vocabulary match that your own Rem.\,2.1 says ``has earned
  nothing''. You have two options. Exhibit the structure map from weld triples to $\Norm$ triples.
  Or retitle the note as being about $\Norm$ (the DLRS sesquilaw magmoid) and demote the weld to
  motivation.
\item \textbf{p.\,3, Claim 4.1, ``Why the receptacle is empty''.} The general sentence ``one cannot
  carry the non-associativity as a cohomology class'' is false. The octonions are the twisted group
  algebra $k_F[\mathbb{Z}_2^3]$ for a 2-cochain $F$ that is \emph{not} a cocycle. Their associator
  $\phi=\partial F$ is a 3-cocycle (Albuquerque--Majid, \emph{Quasialgebra structure of the
  octonions}, J.\ Algebra 220 (1999)). I recomputed this: 168/512 basis triples are non-associative,
  $\partial\phi=1$, the algebra is alternative, and the norm is multiplicative (C9). Quasi-Hopf
  algebras (Drinfeld) are the same picture. In both, the standard cohomological home of
  non-associativity is in degree 3, on an associative base, and that is exactly your setting: $A$ is
  associative and $\#$ is a deformation of it. Moreover $\Norm$ carries a 2-cochain of its own:
  $f\#g=\varphi(f,g)\cdot fg$ with $\varphi(f,g)=D_{fg}$, valued in positive diagonals (see M4). What
  you have actually shown is a type statement: $k$ is a unary 1-cell, so it is not a 2-cocycle. State
  that. Then either test the degree-3 / $\varphi$-valued candidate or drop the general sentence, and
  revisit the registry node. ``Román's magmads'' (l.\,130) is uncited and cannot carry weight here.
\item \textbf{p.\,8, Prop.\,10.1, Cor.\,10.2; abstract (``provably operator-level, never
  element-level''); Rem.\,10.3; \S12.} This is the scope issue of Q2(i)--(iii). Under Def.\,7.1 (U),
  ``principal'' means $K=\mathrm{id}$. What is proved does not cover ``element-level''
  ($\Norm$ itself is $D_x\cdot x$). Replace it with the closure criterion: if $\Im K$ is closed under
  $\cdot$, then $\#=\cdot|_{\Im K}$, so a handed $K$ has a non-closed image. Rewrite the abstract
  clause and the fourth row of the table to match. Typing: in a multi-object $A$, ``a fixed
  $\kappa\in A$'' should be a family $\kappa_X\in A(X,X)$.
\item \textbf{Rem.\,10.3 last line (``its home is original ground''); \S12; Rem.\,11.1.}
  Def.\,7.1 contains the classical transversal construction. Take a group $G$, a subgroup $H$, and a
  right transversal $S\ni 1$. Set $K(g)$ to be the representative of $Hg$. Then (K1), (U) and (LA)
  hold, the last because $K(g)\in Hg$ gives $K(g)c\in Hgc$. The operation $s\#t=K(st)$ is the
  right-loop operation on a transversal (Baer, \emph{Nets and groups}, 1939; R.\ Lal,
  \emph{Transversals in groups}, J.\ Algebra 181 (1996); Kiechle, LNM 1778). In $S_3$ with $|H|=2$,
  3 of the 4 normalized transversals give a non-associative $\#$, and the subgroup $A_3$ gives an
  associative one. With $H$ normal, all are associative (C11). Your structure in general: given (K1),
  (LA) is \emph{equivalent} to $\ker K$ being a right congruence ($K(x)=K(y)\Rightarrow K(xc)=K(yc)$),
  so a normalizer is a right congruence together with a transversal containing the identities.
  $\Norm$ is the multi-object case with $H_X=$ positive diagonals acting on the left: on rows of
  positive mass, every substochastic $x$ factors uniquely as $D\cdot K(x)$. That is an exact
  factorization (Zappa--Szép-type, which is a pleasant irony given the weld). Rem.\,11.1 is honestly
  hedged, but Rem.\,10.3 and \S12 assert originality outright. Cite this literature, delete or scope
  ``original ground'', and re-run the sweep with the keywords \emph{transversal, right loop,
  right congruence, exact factorization}. I have not checked the locators beyond the titles; please
  verify them.
\end{enumerate}

\subsection*{SHOULD}
\begin{enumerate}[label=\textbf{S\arabic*.},leftmargin=*]
\item \textbf{p.\,7, Thm.\,8.2; \S12 ``a theorem rather than a slogan''.} Given Thm.\,8.1, (RA$_\#$)
  is the associativity equation rewritten, so Thm.\,8.2 holds by definition. Call it a corollary or
  an observation; the content is Thm.\,8.1. Also, (RA$_\#$) is not the mirror of (LA): the mirror,
  $K(c\cdot K(x))=K(c\cdot x)$ for all $c$ and $x$, is a different and stronger condition. Calling
  it ``two-sided'' invites the wrong reading.
\item \textbf{p.\,4 Def.\,6.1 vs p.\,5 Prop.\,6.4.} Def.\,6.1 enriches in $\mathbf{Rel}$ under $\subseteq$.
  Prop.\,6.4 switches to the pointwise sub-probability order, which it calls ``the order under which
  the Schoen/Tholen framework applies''; that is asserted, not shown. In the $\subseteq$/support order
  of Def.\,6.1 the two bracketings are \emph{equal}: $f\cdot K(gh)$ and $fgh$ have the same support
  because $D_{gh}$ is a positive diagonal (C3: 0/3000 support mismatches). So in that order the defect
  is invisible, not incomparable. In the pointwise order I confirm incomparability (C3: 0 of 240
  differing triples are comparable). State both; the verdict survives.
\item \textbf{p.\,4, Böhm--Lack--Street paragraph.} Demimonads and weak distributive laws weaken the
  \emph{unit} law, and idempotent splitting does repair unit weakenings (Böhm; Garner's weak
  distributive laws). Your sesquilaw fails \emph{multiplicativity}. So the BLS result is ``Karoubi
  absorbs unit failures'', not ``exactly Corollary 5.4''. Make that contrast the point; it sharpens
  Refutation~II. Drop ``upgrades the verdict''.
\item \textbf{p.\,2 l.\,77--80; \S11 table; email.} The triangulation framing, as answered in Q1.
\item \textbf{p.\,9, Rem.\,10.4.} The multiplication table is incomplete because $b\cdot a$ is not
  given. Both $b\cdot a=z$ and $b\cdot a=a$ give monoids. Only $b\cdot a=z$ satisfies (LA), and with
  $b\cdot a=a$ the witness collapses to $(a\#a)\#a=a=a\#(a\#a)$ (C7). Print $b\cdot a=z$. Add the
  $b\cdot a=a$ table to the Lean file as a breaking control. Add the uniqueness at order 4 (C10).
\item \textbf{p.\,8, Rem.\,9.4.} The recovery regime is ``$h$ total''; $g$ need not be stochastic.
  With $h$ total and 40\% of $g$'s rows equal to $\perp$ I get 0/2000 non-associative triples. The
  control, where $h$ has a $\perp$ row and the middle dimension is at least 3, gives 1085/2000 (C4).
  ``Reproduces DLRS Prop.\,4.9 exactly'' cannot be checked from the note. Quote 4.9's conditions.
\end{enumerate}

\subsection*{NIT}
\begin{enumerate}[label=\textbf{N\arabic*.},leftmargin=*]
\item p.\,3--4, proof of Thm.\,5.2: you write ``$e=s;r$ with $r;s=\mathrm{id}$'' and then set $r=n$,
  $s=m$, which gives $e=m;n=\mathrm{id}_{MD}$. Swap the labels ($s=n$, $r=m$).
\item p.\,3, after Prop.\,3.2: ``rounds 1 and 2'' should read: four rounds, endorsed in round 4 on
  \texttt{40f5c95}, NIT applied in \texttt{b9a0a5c}.
\item p.\,8, ``(see \S10)'' for middle dimension ${\ge}3$: \S10 does not contain this. Give the
  three-line reason instead. $h$ needs a $\perp$ row to create unequal masses, and two
  non-proportional live rows so that the rows of $gh$ are not all parallel. The outer dimensions must
  also be ${\ge}2$. So $1\to2\to3\to2$ is minimal (C5: shapes with middle 2, or an outer 1, give 0/600).
\item p.\,5, near-miss: $\mathrm{id}_{DM}\le k$ holds pointwise on the $Y$-coordinates only. At the
  $\perp$-coordinate of $D(1+Y)$ the inequality reverses.
\item p.\,8, Cor.\,10.2: ``linear'' has no meaning in an abstract $A$. Use ``$K\ne\mathrm{id}$''
  together with (U) (see M3).
\item ``magmad'' and ``magmoid'' are used interchangeably (pp.\,3--4). Pick one.
\item p.\,3, Rem.\,3.3: name the registry node for $[\omega]\in H^2$.
\end{enumerate}

\section*{Checks run (script output in the \texttt{.log})}
{\small
\begin{tabular}{lll}
\toprule
C1 & Lemma 9.3 witness, exact & matches PDF \\
C2 & (LA), row-$K$, 3000 random; \textbf{control} column-$K$ & 0 / 945 violations\\
C3 & Thm 8.1; pointwise comparability; support equality & OK; 0/240; 0 mismatches\\
C4 & $h$ total $\Rightarrow$ assoc.; \textbf{control} $h$ with $\perp$ row & 0/2000; 1085/2000\\
C5 & minimal shape & only $(1,2,{\ge}3,2)$ has defects\\
C6 & handed normalizers in all monoids, order 1--4 & 0,0,0,6; all with non-closed $\Im K$\\
C7 & Lean model, $b\cdot a\in\{z,a\}$ (\textbf{control}) & only $z$ works\\
C8 & principal $\kappa x$, $x\kappa$, $\kappa x\kappa$, no (U) & 1434 tested, 0 non-assoc; (U) $\Rightarrow$ id\\
C9 & octonions $=k_F[\mathbb{Z}_2^3]$, $\phi=\partial F$ & 3-cocycle, 168/512 non-assoc\\
C10 & order-4 handed models up to isomorphism & 1 (your model, $b\cdot a=z$)\\
C11 & $S_3$ transversals as normalizers & (K1),(U),(LA) all hold; 3/4 non-assoc\\
\bottomrule
\end{tabular}}

Not checked: DLRS Prop.\,4.9 itself, Schoen/Tholen beyond your Def.\,6.1, the Lean source (only your
description of it), and the BLS demimonad locator.
\end{document}
